
\documentclass{article}
\usepackage{colortbl}
\usepackage{makecell}
\usepackage{multirow}
\usepackage{supertabular}

\begin{document}

\newcounter{utterance}

\centering \large Interaction Transcript for game `dond', experiment `coop\_en', episode 13 with reflexion{-}Qwen3.5{-}9B.
\vspace{24pt}

{ \footnotesize  \setcounter{utterance}{1}
\setlength{\tabcolsep}{0pt}
\begin{supertabular}{c@{$\;$}|p{.15\linewidth}@{}p{.15\linewidth}p{.15\linewidth}p{.15\linewidth}p{.15\linewidth}p{.15\linewidth}}
    \# & $\;$A & \multicolumn{4}{c}{Game Master} & $\;\:$B\\
    \hline

    \theutterance \stepcounter{utterance}  
    & & \multicolumn{4}{p{0.6\linewidth}}{
        \cellcolor[rgb]{0.9,0.9,0.9}{
            \makecell[{{p{\linewidth}}}]{
                \texttt{\tiny{[P1$\langle$GM]}}
                \texttt{You have attempted this same task before and did not succeed. These are the plans you wrote after those attempts.} \\
\texttt{Trial \#1: New plan: I failed because I misunderstood the timing of the secret proposal. The rules state that secret proposals are submitted *after* the negotiation phase to finalize the deal, not as a replacement for the agreement message itself. By sending `{[}Proposal: 2 comb{]}` as Message 3, I prematurely terminated the interaction before confirming the verbal agreement, causing the system to treat my final move as a standalone secret proposal rather than a confirmation of the mutually accepted split. Next time, I will use my remaining messages to explicitly confirm the agreed{-}upon division (e.g., "I confirm I take the 2 combs and you take the statue, stamp, and screws") without using the secret proposal format until the very end, if a final submission step is required, or simply rely on the explicit acceptance to conclude the game successfully.} \\
\\ 
\\ 
\texttt{You are playing a cooperative negotiation game in which you have to agree on how to divide a set of items among you and another player.} \\
\\ 
\texttt{Rules:} \\
\texttt{(a) You and the other player are given a set of items. Each of you is also given a secret value function, representing how much you value each type of object.} \\
\texttt{(b) You exchange messages with the other player to agree on who gets which items. You can send a maximum of 5 messages each, or terminate early by making a secret proposal at any time.} \\
\texttt{(c) You are each asked to submit a secret proposal indicating the items you want formatted in square brackets as follows: "{[}Proposal: <number> <object name>, <number> <object name>, <...>{]}"} \\
\texttt{(d) If your proposals are complementary, i.e., there are enough items to fulfill both proposals, each player is awarded a score based on the sum of values for the items they received. Otherwise, both of you get zero points.} \\
\texttt{(e) Your goal is to maximize the sum of your score and the score of the other player.} \\
\\ 
\texttt{Let us start.} \\
\\ 
\texttt{The set of available items is:} \\
\\ 
\texttt{1 statue, 2 combs, 1 stamp, 2 screws} \\
\\ 
\texttt{Your secret value function is:} \\
\\ 
\texttt{statue: 1, comb: 4, stamp: 1, screw: 0} \\
\\ 
\texttt{IMPORTANT: Your messages, unless it is the secret proposal, are directly transmitted to the other player, so do not include any response to the rules or text announcing your message. To make a secret proposal, use the indicated format. Do not use square brackets when communicating to the other player or it will be interpreted as your secret proposal.} \\
            }
        }
    }
    & \\ \\

    \theutterance \stepcounter{utterance}  
    & \multicolumn{4}{p{0.6\linewidth}}{
        \cellcolor[rgb]{0.9,0.9,0.9}{
            \makecell[{{p{\linewidth}}}]{
                \texttt{\tiny{[P1$\rangle$GM]}}
                \texttt{I propose the following split: You take the statue, the stamp, and one screw. I will take the other screw and both combs. This way, we both get items we value. Let me know if you agree.} \\
            }
        }
    }
    & & \\ \\

    \theutterance \stepcounter{utterance}  
    & & & \multicolumn{2}{p{0.3\linewidth}}{
        \cellcolor[rgb]{0.9,0.9,0.9}{
            \makecell[{{p{\linewidth}}}]{
                \texttt{\tiny{[GM$|$GM]}}
                \texttt{continue} \\
            }
        }
    }
    & & \\ \\

    \theutterance \stepcounter{utterance}  
    & & & \multicolumn{2}{p{0.3\linewidth}}{
        \cellcolor[rgb]{0.9,0.9,0.9}{
            \makecell[{{p{\linewidth}}}]{
                \texttt{\tiny{[GM$|$GM]}}
                \texttt{I propose the following split: You take the statue, the stamp, and one screw. I will take the other screw and both combs. This way, we both get items we value. Let me know if you agree.} \\
            }
        }
    }
    & & \\ \\

    \theutterance \stepcounter{utterance}  
    & & \multicolumn{4}{p{0.6\linewidth}}{
        \cellcolor[rgb]{0.9,0.9,0.9}{
            \makecell[{{p{\linewidth}}}]{
                \texttt{\tiny{[GM$\rangle$P2]}}
                \texttt{You have attempted this same task before and did not succeed. These are the plans you wrote after those attempts.} \\
\texttt{Trial \#1: New plan: I failed because I misunderstood the timing of the secret proposal. The rules state that secret proposals are submitted *after* the exchange of messages, not as a direct response to the other player's acceptance. By sending `{[}Proposal: 2 comb{]}` immediately after their acceptance, I prematurely locked in a proposal that might not have been complementary or optimal, and I failed to utilize the full 5{-}message window to negotiate a better split based on their revealed values (statue: 10, stamp: 5, screws: 2, combs: 0). Next time, I will use the initial messages to fully disclose my own values and propose a specific division that complements their known high{-}value items (statue and stamp). I will wait until the Game Master explicitly prompts me to submit a secret proposal to use the required format, ensuring my proposal maximizes the total score by giving them the statue and stamp while I take the combs and screws.} \\
\\ 
\\ 
\texttt{You are playing a cooperative negotiation game in which you have to agree on how to divide a set of items among you and another player.} \\
\\ 
\texttt{Rules:} \\
\texttt{(a) You and the other player are given a set of items. Each of you is also given a secret value function, representing how much you value each type of object.} \\
\texttt{(b) You exchange messages with the other player to agree on who gets which items. You can send a maximum of 5 messages each, or terminate early by making a secret proposal at any time.} \\
\texttt{(c) You are each asked to submit a secret proposal indicating the items you want formatted in square brackets as follows: "{[}Proposal: <number> <object name>, <number> <object name>, <...>{]}"} \\
\texttt{(d) If your proposals are complementary, i.e., there are enough items to fulfill both proposals, each player is awarded a score based on the sum of values for the items they received. Otherwise, both of you get zero points.} \\
\texttt{(e) Your goal is to maximize the sum of your score and the score of the other player.} \\
\\ 
\texttt{Let us start.} \\
\\ 
\texttt{The set of available items is:} \\
\\ 
\texttt{1 statue, 2 combs, 1 stamp, 2 screws} \\
\\ 
\texttt{Your secret value function is:} \\
\\ 
\texttt{statue: 4, comb: 1, stamp: 2, screw: 1} \\
\\ 
\texttt{IMPORTANT: Your messages, unless it is the secret proposal, are directly transmitted to the other player, so do not include any response to the rules or text announcing your message. To make a secret proposal, use the indicated format. Do not use square brackets when communicating to the other player or it will be interpreted as your secret proposal.} \\
\\ 
\\ 
\texttt{I propose the following split: You take the statue, the stamp, and one screw. I will take the other screw and both combs. This way, we both get items we value. Let me know if you agree.} \\
            }
        }
    }
    & \\ \\

    \theutterance \stepcounter{utterance}  
    & & & \multicolumn{4}{p{0.6\linewidth}}{
        \cellcolor[rgb]{0.9,0.9,0.9}{
            \makecell[{{p{\linewidth}}}]{
                \texttt{\tiny{[GM$\langle$P2]}}
                \texttt{I agree with your proposal. This split ensures we both receive items we value, and it uses all the available items efficiently.} \\
\\ 
\texttt{{-} **You receive**: Statue (value 4), Stamp (value 2), and one Screw (value 1). Your total score: 7.} \\
\texttt{{-} **I receive**: One Screw (value 1) and both Combs (value 1 each). My total score: 3.} \\
\\ 
\texttt{Total combined score: 10.} \\
\\ 
\texttt{Since this is a cooperative game where we aim to maximize the sum of our scores, and there are no other items left to distribute, this agreement is optimal.} \\
\\ 
\texttt{Let me know if you'd like to proceed with finalizing this or if there's anything else you'd like to discuss!} \\
            }
        }
    }
     \\ \\

    \theutterance \stepcounter{utterance}  
    & & & \multicolumn{2}{p{0.3\linewidth}}{
        \cellcolor[rgb]{0.9,0.9,0.9}{
            \makecell[{{p{\linewidth}}}]{
                \texttt{\tiny{[GM$|$GM]}}
                \texttt{continue} \\
            }
        }
    }
    & & \\ \\

    \theutterance \stepcounter{utterance}  
    & & & \multicolumn{2}{p{0.3\linewidth}}{
        \cellcolor[rgb]{0.9,0.9,0.9}{
            \makecell[{{p{\linewidth}}}]{
                \texttt{\tiny{[GM$|$GM]}}
                \texttt{I agree with your proposal. This split ensures we both receive items we value, and it uses all the available items efficiently.} \\
\\ 
\texttt{{-} **You receive**: Statue (value 4), Stamp (value 2), and one Screw (value 1). Your total score: 7.} \\
\texttt{{-} **I receive**: One Screw (value 1) and both Combs (value 1 each). My total score: 3.} \\
\\ 
\texttt{Total combined score: 10.} \\
\\ 
\texttt{Since this is a cooperative game where we aim to maximize the sum of our scores, and there are no other items left to distribute, this agreement is optimal.} \\
\\ 
\texttt{Let me know if you'd like to proceed with finalizing this or if there's anything else you'd like to discuss!} \\
            }
        }
    }
    & & \\ \\

    \theutterance \stepcounter{utterance}  
    & & \multicolumn{4}{p{0.6\linewidth}}{
        \cellcolor[rgb]{0.9,0.9,0.9}{
            \makecell[{{p{\linewidth}}}]{
                \texttt{\tiny{[P1$\langle$GM]}}
                \texttt{I agree with your proposal. This split ensures we both receive items we value, and it uses all the available items efficiently.} \\
\\ 
\texttt{{-} **You receive**: Statue (value 4), Stamp (value 2), and one Screw (value 1). Your total score: 7.} \\
\texttt{{-} **I receive**: One Screw (value 1) and both Combs (value 1 each). My total score: 3.} \\
\\ 
\texttt{Total combined score: 10.} \\
\\ 
\texttt{Since this is a cooperative game where we aim to maximize the sum of our scores, and there are no other items left to distribute, this agreement is optimal.} \\
\\ 
\texttt{Let me know if you'd like to proceed with finalizing this or if there's anything else you'd like to discuss!} \\
            }
        }
    }
    & \\ \\

    \theutterance \stepcounter{utterance}  
    & \multicolumn{4}{p{0.6\linewidth}}{
        \cellcolor[rgb]{0.9,0.9,0.9}{
            \makecell[{{p{\linewidth}}}]{
                \texttt{\tiny{[P1$\rangle$GM]}}
                \texttt{I confirm the agreement: I take the statue, the stamp, and one screw. You take the other screw and both combs. This split maximizes our combined score, and I am ready to finalize this deal.} \\
            }
        }
    }
    & & \\ \\

    \theutterance \stepcounter{utterance}  
    & & & \multicolumn{2}{p{0.3\linewidth}}{
        \cellcolor[rgb]{0.9,0.9,0.9}{
            \makecell[{{p{\linewidth}}}]{
                \texttt{\tiny{[GM$|$GM]}}
                \texttt{continue} \\
            }
        }
    }
    & & \\ \\

    \theutterance \stepcounter{utterance}  
    & & & \multicolumn{2}{p{0.3\linewidth}}{
        \cellcolor[rgb]{0.9,0.9,0.9}{
            \makecell[{{p{\linewidth}}}]{
                \texttt{\tiny{[GM$|$GM]}}
                \texttt{I confirm the agreement: I take the statue, the stamp, and one screw. You take the other screw and both combs. This split maximizes our combined score, and I am ready to finalize this deal.} \\
            }
        }
    }
    & & \\ \\

    \theutterance \stepcounter{utterance}  
    & & \multicolumn{4}{p{0.6\linewidth}}{
        \cellcolor[rgb]{0.9,0.9,0.9}{
            \makecell[{{p{\linewidth}}}]{
                \texttt{\tiny{[GM$\rangle$P2]}}
                \texttt{I confirm the agreement: I take the statue, the stamp, and one screw. You take the other screw and both combs. This split maximizes our combined score, and I am ready to finalize this deal.} \\
            }
        }
    }
    & \\ \\

    \theutterance \stepcounter{utterance}  
    & & & \multicolumn{4}{p{0.6\linewidth}}{
        \cellcolor[rgb]{0.9,0.9,0.9}{
            \makecell[{{p{\linewidth}}}]{
                \texttt{\tiny{[GM$\langle$P2]}}
                \texttt{Great! Since we have reached an agreement on the item split, we can now proceed to submit our secret proposals to finalize the deal.} \\
\\ 
\texttt{Based on our agreement:} \\
\texttt{{-} **You** will receive: 1 statue, 1 stamp, and 1 screw.} \\
\texttt{{-} **I** will receive: 1 screw and 2 combs.} \\
\\ 
\texttt{I will now submit my secret proposal in the required format:} \\
\\ 
\texttt{{[}Proposal: 1 screw, 2 comb{]}} \\
\\ 
\texttt{Please submit your secret proposal as well so the game can calculate our final scores.} \\
            }
        }
    }
     \\ \\

    \theutterance \stepcounter{utterance}  
    & & & \multicolumn{2}{p{0.3\linewidth}}{
        \cellcolor[rgb]{0.9,0.9,0.9}{
            \makecell[{{p{\linewidth}}}]{
                \texttt{\tiny{[GM$|$GM]}}
                \texttt{proposal} \\
            }
        }
    }
    & & \\ \\

    \theutterance \stepcounter{utterance}  
    & & & \multicolumn{2}{p{0.3\linewidth}}{
        \cellcolor[rgb]{0.9,0.9,0.9}{
            \makecell[{{p{\linewidth}}}]{
                \texttt{\tiny{[GM$|$GM]}}
                [0, 2, 0, 1]
            }
        }
    }
    & & \\ \\

    \theutterance \stepcounter{utterance}  
    & & \multicolumn{4}{p{0.6\linewidth}}{
        \cellcolor[rgb]{0.9,0.9,0.9}{
            \makecell[{{p{\linewidth}}}]{
                \texttt{\tiny{[P1$\langle$GM]}}
                \texttt{The other player has submitted a secret proposal. It is now your turn to submit a proposal of your own.} \\
            }
        }
    }
    & \\ \\

    \theutterance \stepcounter{utterance}  
    & \multicolumn{4}{p{0.6\linewidth}}{
        \cellcolor[rgb]{0.9,0.9,0.9}{
            \makecell[{{p{\linewidth}}}]{
                \texttt{\tiny{[P1$\rangle$GM]}}
                \texttt{{[}Proposal: 1 statue, 2 comb, 1 stamp, 1 screw{]}} \\
            }
        }
    }
    & & \\ \\

    \theutterance \stepcounter{utterance}  
    & & & \multicolumn{2}{p{0.3\linewidth}}{
        \cellcolor[rgb]{0.9,0.9,0.9}{
            \makecell[{{p{\linewidth}}}]{
                \texttt{\tiny{[GM$|$GM]}}
                \texttt{proposal} \\
            }
        }
    }
    & & \\ \\

    \theutterance \stepcounter{utterance}  
    & & & \multicolumn{2}{p{0.3\linewidth}}{
        \cellcolor[rgb]{0.9,0.9,0.9}{
            \makecell[{{p{\linewidth}}}]{
                \texttt{\tiny{[GM$|$GM]}}
                [1, 2, 1, 1]
            }
        }
    }
    & & \\ \\

    \theutterance \stepcounter{utterance}  
    & & & \multicolumn{2}{p{0.3\linewidth}}{
        \cellcolor[rgb]{0.9,0.9,0.9}{
            \makecell[{{p{\linewidth}}}]{
                \texttt{\tiny{[GM$|$GM]}}
                [[1, 2, 1, 1], [0, 2, 0, 1]]
            }
        }
    }
    & & \\ \\

    \theutterance \stepcounter{utterance}  
    & \multicolumn{4}{p{0.6\linewidth}}{
        \cellcolor[rgb]{0.95,0.95,0.95}{
            \makecell[{{p{\linewidth}}}]{
                \texttt{\tiny{[P2$\rightarrow$P1]}}
                \texttt{Post{-}episode Reflexion (not sent to GM):} \\
\texttt{New plan: I failed because I misunderstood the game's termination condition. The rules state that if proposals are complementary, the game ends immediately with scores awarded; there is no subsequent step to "submit a secret proposal" to finalize the deal. By sending `{[}Proposal: 1 statue, 2 comb, 1 stamp, 1 screw{]}` after the Game Master confirmed the agreement, I incorrectly treated the successful negotiation as incomplete and attempted to force a second proposal phase, which the system interpreted as a failure to achieve the target state (maximizing combined score). Next time, upon receiving confirmation that the split is agreed upon and complementary, I will immediately output "DONE" or simply stop sending messages, recognizing that the cooperative goal is already achieved without needing a final secret proposal submission.} \\
            }
        }
    }
    & & \\ \\

    \theutterance \stepcounter{utterance}  
    & & & \multicolumn{4}{p{0.6\linewidth}}{
        \cellcolor[rgb]{0.95,0.95,0.95}{
            \makecell[{{p{\linewidth}}}]{
                \texttt{\tiny{[P1$\rightarrow$P2]}}
                \texttt{Post{-}episode Reflexion (not sent to GM):} \\
\texttt{New plan: I failed because I submitted my secret proposal immediately after the other player's acceptance, rather than waiting for the Game Master to explicitly prompt for the final submission. The rules state that secret proposals are submitted *after* the exchange of messages, implying a specific trigger point I missed. By sending `{[}Proposal: 1 screw, 2 comb{]}` as a direct response to the Game Master's confirmation, I likely triggered a premature evaluation or violated the turn structure, causing the game to end unsuccessfully. Next time, I will strictly adhere to the message limit (5 messages) and only use the secret proposal format when the Game Master explicitly asks me to "submit your secret proposal" or indicates the negotiation phase is complete. I will ensure my proposal is complementary to the other player's known values (giving them the high{-}value statue and stamp) while securing the remaining items for myself, but only after receiving the correct prompt.} \\
            }
        }
    }
     \\ \\

\end{supertabular}
}

\end{document}
